\section{Conformally weighted chords}\label{sec:chords}

Throughout this section, $n>2$ and the normalized extremizing cycle
$A$ is as in Section~\ref{sec:radial}. We assume the sharp filling
inequality in dimension $n-1$, so that Sections~\ref{sec:sobolev}
and~\ref{sec:minimizer} apply. Keep the notation
\[
 \beta=\frac1{n-2},\qquad p=\frac{2n}{n-2},\qquad
 Q(u)=4\beta E(u)+n\norm{u}_2^2.
\]
Let $v\ge0$ be the quotient minimizer from
Proposition~\ref{prop:minimizer}. In particular,
\begin{equation}\label{chord:normalization}
 0<W:=\int_Y v^p\dd\mu<s_n,\qquad Q(v)=nW,
 \qquad \dd\nu=\frac{v^p}{W}\dd\mu.
\end{equation}
The last expression defines a probability measure. By
Lemma~\ref{lem:powers}, $v$ has all finite positive moments and
\begin{equation}\label{chord:weighted-energy}
 \int_Y v^q\abs{Dv}^2\dd\mu<\infty\qquad(q\ge0).
\end{equation}
Choose a finite nonnegative Borel representative of $v$; changing
it on a $\mu$-null set has no effect on these assertions.

For a fixed center $y\in Y$ with $v(y)>0$, set
\begin{equation}\label{chord:data}
 \kappa=v(y)^\beta,\qquad r(x)=d(y,x),\qquad
 s(x)=\kappa r(x)v(x)^\beta,
 \qquad L_y=\int_Y s^2\dd\nu.
\end{equation}
These are finite Borel functions or numbers. In particular, with
$D=\diam Y$,
\[
 L_y\le\frac{\kappa^2D^2}{W}
             \int_Y v^{p+2\beta}\dd\mu<\infty.
\]
We study the chord distribution through the transform
\begin{equation}\label{chord:transform-definition}
 J_y(a)=\int_Y(1+as^2)^{-n}\dd\nu\qquad(a\ge0).
\end{equation}
Our first comparison will give
$J_y(a)\le(1+2aL_y)^{-n/2}$ for every center with $v(y)>0$.
Euclidean tangent density will then give a lower bound for
$\liminf_{a\to\infty}a^{n/2}J_y(a)$ at almost every center, forcing
\[
 L_y\le2(W/s_n)^{2/n}<2
 \quad\text{for $\nu$-almost every center $y$}.
\]
Section~\ref{sec:conclusion} will show that the average of these same
moments is at least $2$.

Every chart derivative in this section is taken in the variable
$x$, at this fixed center $y$. We shall not require $s\in\Hcal$.
Instead, define the square-integrable chart covector field
\begin{equation}\label{chord:xi}
 \xi=\kappa\bigl(v^{1+\beta}Dr
                   +\beta r v^\beta Dv\bigr).
\end{equation}
Its square integrability follows from $\abs{Dr}\le1$, boundedness
of $r$, the moments of $v$, and~\eqref{chord:weighted-energy} with
$q=2\beta$. On $\{v>0\}$ it is the formal expression $vDs$;
the explicit formula~\eqref{chord:xi} defines it also on $\{v=0\}$.

\subsection{Admissible chord tests}

For $a>0$, the profile $w(t)=(1+at^2)^{-(n-2)/2}$ is the rescaled
Euclidean Sobolev profile used in Lemma~\ref{lem:radial-sobolev}.
Its exponent is suited to the transform because
$w(s)^p=(1+as^2)^{-n}$. We will compare the quotient of $vw(s)$
with that of $v$, using $vw(s)^2$ in the minimizer equation.
The following lemma first justifies these tests and the radial tests
needed to control their energy, including at points where $v=0$.

\begin{lemma}[Admissibility of the chord tests]\label{lem:chord-tests}
Fix $a>0$ and put
\[
 w(t)=(1+at^2)^{-(n-2)/2}\qquad(t\in\R).
\]
Then $vw(s)$ and $vw(s)^2$ belong to $\Hcal$, and
\begin{equation}\label{chord:product-rules}
 \begin{split}
 D(vw(s))&=w(s)Dv+w'(s)\xi,\\
 D(vw(s)^2)&=w(s)^2Dv+2w(s)w'(s)\xi.
 \end{split}
\end{equation}
If $f:\R\to[0,\infty)$ is smooth with bounded value and bounded
first derivative, then
\begin{equation}\label{chord:radial-tests}
 g=\kappa^2v^{2+2\beta}f(s),\qquad
 \psi=\beta\kappa^2r^2v^{1+2\beta}f(s)
\end{equation}
also belong to $\Hcal$. In particular, $g$ is an admissible
nonnegative test in the extended radial inequality, and $\psi$
is an admissible test in~\eqref{eq:euler}.
\end{lemma}

\begin{proof}
We verify the hypotheses of Lemma~\ref{lem:composites} for the
actual functions of $(r,v)$ that occur here. This is necessary
when $\beta<1$, since the scalar function $v^\beta$ itself is
not $C^1$ at zero.

For $j=1,2$ and $z>0$, define
\[
 F_j(r,z)=z\,w(\kappa r z^\beta)^j,
 \qquad t=\kappa r z^\beta.
\]
On this half-plane,
\[
 \partial_zF_j=w(t)^j+j\beta t\,w(t)^{j-1}w'(t),\qquad
 \partial_rF_j=j\kappa z^{1+\beta}w(t)^{j-1}w'(t).
\]
As $z\downarrow0$, uniformly for $r$ in bounded intervals,
$w(t)=1+O(t^2)$ and $w'(t)=O(t)$. Thus
\[
 F_j(r,z)=z+O(z^{1+2\beta}),\qquad
 \partial_zF_j(r,z)\longrightarrow1,
 \qquad \partial_rF_j(r,z)\longrightarrow0.
\]
Extension by $F_j(r,z)=z$ on $z\le0$ is therefore jointly $C^1$.
The functions and their first derivatives have polynomial growth
in $z\ge0$ on bounded $r$ intervals.

For the other tests, write
\[
 F_g(r,z)=\kappa^2z^{2+2\beta}f(\kappa r z^\beta),\qquad
 F_\psi(r,z)=\beta\kappa^2r^2z^{1+2\beta}
                          f(\kappa r z^\beta)
 \quad(z>0).
\]
Their partial derivatives are
\begin{align*}
 \partial_zF_g
   &=\kappa^2z^{1+2\beta}
                  \bigl((2+2\beta)f(t)+\beta t f'(t)\bigr),\\
 \partial_rF_g
   &=\kappa^3z^{2+3\beta}f'(t),\\
 \partial_zF_\psi
   &=\beta\kappa^2r^2
       \bigl((1+2\beta)z^{2\beta}f(t)
                    +\beta\kappa r z^{3\beta}f'(t)\bigr),\\
 \partial_rF_\psi
   &=\beta\kappa^2
       \bigl(2r z^{1+2\beta}f(t)
                    +\kappa r^2z^{1+3\beta}f'(t)\bigr).
\end{align*}
Their values and all these derivatives tend to zero as
$z\downarrow0$, uniformly for bounded $r$. The smallest power
of $z$ in the derivative formulas is $2\beta>0$. Consequently,
extension by zero on $z\le0$ makes both functions jointly $C^1$,
even if $f(0)$ or $f'(0)$ is nonzero. Their values and first
derivatives again have polynomial growth on the range of $r$.

These $C^1$ extensions and polynomial bounds satisfy all the
hypotheses of Lemma~\ref{lem:composites}. That lemma gives membership
in $\Hcal$ and the chain rules without a boundedness assumption on $v$.
On $\{v>0\}$, differentiation gives~\eqref{chord:product-rules}.
On $\{v=0\}$, Proposition~\ref{prop:sobolev-space} gives $Dv=0$
almost everywhere, while~\eqref{chord:xi} gives $\xi=0$.
Thus these formulas hold $\mu$-almost everywhere on all of $Y$.
\end{proof}

\begin{lemma}[A consequence of the radial inequality]
\label{lem:completed-radial}
For every center $y\in Y$, with $r=d(y,\cdot)$, and every
nonnegative $g\in\Hcal$,
\begin{equation}\label{chord:completed-radial}
 \int_Y\bigl(ng\abs{Dr}^2+r\ip{Dr}{Dg}\bigr)\dd\mu
       \le\frac n4\int_Y r^2g\dd\mu.
\end{equation}
\end{lemma}

\begin{proof}
The extension of~\eqref{eq:radial} to these tests is part of
Proposition~\ref{prop:sobolev-space}. Pointwise, with
$b=\sqrt{1-\abs{Dr}^2}$,
\[
 1-rb-\abs{Dr}^2+\frac{r^2}{4}
                   =\left(b-\frac r2\right)^2\ge0.
\]
Multiply by $ng\ge0$ and combine this inequality
with~\eqref{eq:radial}. All terms are integrable by boundedness
of $r$, $\abs{Dr}\le1$, finite $\mu$, and $g\in\Hcal$.
\end{proof}

\subsection{Distribution comparison}

All required composite tests are now admissible, including on the zero
set of $v$. We combine quotient minimality with radial variation to
control the chord transform. Once the integral estimate is obtained,
the remaining comparison is a scalar differential inequality.

\begin{proposition}[Distribution comparison]\label{prop:chord-transform}
For every center $y$ with $v(y)>0$, and every $a\ge0$,
\begin{equation}\label{eq:chord-transform}
 J_y(a)=\int_Y(1+as^2)^{-n}\dd\nu
                  \le(1+2aL_y)^{-n/2},
\end{equation}
where $s$ and $L_y$ are defined in~\eqref{chord:data}.
\end{proposition}

\begin{proof}
Fix the center and suppress its subscript on $J$. For $a>0$ let
$w(t)=(1+at^2)^{-(n-2)/2}$ as in
Lemma~\ref{lem:chord-tests}. Since $p(n-2)/2=n$,
\[
 \int_Y w(s)^p\dd\nu=J(a).
\]
Testing~\eqref{eq:euler} with $vw(s)^2$ and using
the two product formulas~\eqref{chord:product-rules} gives
\begin{equation}\label{chord:ground-state}
 Q(vw(s))
    =n\int_Y v^p w(s)^2\dd\mu
              +4\beta\int_Y w'(s)^2\abs{\xi}^2\dd\mu.
\end{equation}
Indeed the terms $w^2\abs{Dv}^2+2ww'\ip{Dv}{\xi}$ in
$\abs{D(vw)}^2$ equal $\ip{Dv}{D(vw^2)}$. The quotient
minimality of $v$ gives, on the other hand,
\begin{equation}\label{chord:quotient}
 Q(vw(s))\ge
   \frac{Q(v)}{\norm{v}_p^2}\norm{vw(s)}_p^2
                   =nW J(a)^{(n-2)/n}.
\end{equation}
The test $vw(s)$ is nonzero because $w$ is strictly positive
and $W>0$.

We next control the last integral
in~\eqref{chord:ground-state}. Let $f$ be smooth, bounded and
nonnegative with bounded derivative, and use $g,\psi$ from
\eqref{chord:radial-tests}. Add
\eqref{chord:completed-radial} to~\eqref{eq:euler} divided by
$4\beta$, with test $\psi$. The identity
$v\psi=\beta r^2g$ cancels the two undifferentiated terms and
yields
\begin{equation}\label{chord:added-tests}
 \int_Y\bigl(ng\abs{Dr}^2+r\ip{Dr}{Dg}
                              +\ip{Dv}{D\psi}\bigr)\dd\mu
       \le\frac n4\int_Y s^2f(s)v^p\dd\mu.
\end{equation}
Here we used
$v^{p-1}\psi=\beta s^2f(s)v^p$.
For completeness, after removing the common factor $\kappa^2$,
the coefficients of $\abs{Dr}^2$, $\ip{Dr}{Dv}$, and
$\abs{Dv}^2$ in the left integrand are respectively
\begin{align*}
 &v^{2+2\beta}\bigl(nf+sf'\bigr),\\
 &r v^{1+2\beta}\bigl((2+4\beta)f+2\beta sf'\bigr),\\
 &\beta r^2v^{2\beta}\bigl((1+2\beta)f+\beta sf'\bigr).
\end{align*}
The scalar functions $f,f'$ in this display are evaluated at
$s$. Since
\begin{equation}\label{chord:exponent-identities}
 1+2\beta=n\beta,\qquad 2+4\beta=2n\beta=p,
\end{equation}
these are precisely the coefficients of
$(nf(s)+sf'(s))\abs{\xi}^2/\kappa^2$. Thus
\begin{equation}\label{chord:energy}
 \int_Y\bigl(nf(s)+sf'(s)\bigr)\abs{\xi}^2\dd\mu
               \le\frac n4\int_Ys^2f(s)v^p\dd\mu.
\end{equation}

To bound the energy term in~\eqref{chord:ground-state}, we want
$nf(t)+tf'(t)=w'(t)^2$ in~\eqref{chord:energy}. Integrating this
first-order equation leads to the choice
\begin{equation}\label{chord:f}
 f(t)=\int_0^1\tau^{n-1}w'(\tau t)^2\dd\tau
                   \qquad(t\in\R).
\end{equation}
This is smooth and nonnegative. Both $w'$ and $w''$ are bounded
on $\R$, so $f$ and
$f'(t)=2\int_0^1\tau^nw'(\tau t)w''(\tau t)\dd\tau$
are bounded. Integrating the derivative of
$\tau^nw'(\tau t)^2$ gives
\begin{equation}\label{chord:f-identity}
 nf(t)+tf'(t)=w'(t)^2.
\end{equation}
The boundary term at $\tau=0$ is zero. Combining
\eqref{chord:ground-state}, \eqref{chord:quotient},
\eqref{chord:energy}, and~\eqref{chord:f-identity}, and dividing
by $nW$, we obtain
\begin{equation}\label{chord:scalar-start}
 J(a)^{(n-2)/n}
          \le\int_Y\bigl(w(s)^2+\beta s^2f(s)\bigr)\dd\nu.
\end{equation}

The current and variational estimates have now reduced the comparison
to \eqref{chord:scalar-start}. We evaluate its scalar right side and
solve the resulting differential inequality for the transform. Put
$b=as^2\ge0$ and $q=(n-2)/2>0$. Since
\[
 w'(t)^2=(n-2)^2a^2t^2(1+at^2)^{-n},
\]
the substitution $u=\tau^2$ in~\eqref{chord:f} gives
\[
 w(s)^2+\beta s^2f(s)
 =(1+b)^{-(n-2)}
             +q b^2\int_0^1\frac{u^{n/2}}{(1+bu)^n}\dd u.
\]
The exact derivative identity
\[
 \frac{\dd}{\dd u}\left(\frac{u^q}{(1+bu)^{n-2}}\right)
       =q\frac{u^{q-1}(1-b^2u^2)}{(1+bu)^n}
\]
has an integrable right side on $(0,1)$; this includes $n=3$,
when $q-1=-1/2$. Integrating it shows that
\begin{equation}\label{chord:scalar-identity}
 w(s)^2+\beta s^2f(s)
       =q\int_0^1\frac{u^{q-1}}{(1+as^2u)^n}\dd u.
\end{equation}
In particular this expression lies in $(0,1]$. Tonelli's theorem
therefore turns~\eqref{chord:scalar-start} into
\begin{equation}\label{chord:J-G}
 J(a)^{(n-2)/n}\le G(a),\qquad
 G(a)=q\int_0^1J(au)u^{q-1}\dd u.
\end{equation}

We have $G(0)=1$ and $G(a)>0$. For $a>0$, changing variable
$t=au$ gives
\[
 G(a)=q a^{-q}\int_0^a J(t)t^{q-1}\dd t,
 \qquad aG'(a)=q\bigl(J(a)-G(a)\bigr).
\]
Continuity of $J$ suffices for this differentiation. The finite
second moment in~\eqref{chord:data} also allows differentiation
at zero from the right, because
\[
 \left|\frac{\partial}{\partial a}(1+as^2)^{-n}\right|
                 \le ns^2\qquad(a\ge0).
\]
Consequently
\begin{equation}\label{chord:initial-derivatives}
 J'(0+)=-nL_y,
 \qquad G'(0+)=\frac{q}{q+1}J'(0+)=-(n-2)L_y.
\end{equation}
Set $h=G^{-1/q}=G^{-2/(n-2)}$. For $a>0$,
\[
 ah'=h-JG^{-n/(n-2)}\ge h-1,
\]
where the last step follows from~\eqref{chord:J-G}. Thus
\[
 \left(\frac{h(a)-1}{a}\right)'
       =\frac{ah'(a)-h(a)+1}{a^2}\ge0.
\]
By~\eqref{chord:initial-derivatives}, the quotient on the left
has limit $2L_y$ at zero. It follows that
$h(a)\ge1+2aL_y$. Finally~\eqref{chord:J-G} gives
\[
 J(a)\le G(a)^{n/(n-2)}=h(a)^{-n/2}
                          \le(1+2aL_y)^{-n/2}.
\]
At $a=0$ both sides equal one. If $L_y=0$, then $s=0$
$\nu$-almost everywhere, so $J=G=h=1$ and the same conclusion
holds. No higher chord moment is needed in the scalar comparison.
\end{proof}

\subsection{Tangent density}

The transform estimate bounds $J_y$ in terms of its second moment
$L_y$, but does not yet bound that moment. At almost every center,
integer Euclidean tangent density supplies a lower coefficient for
$J_y$ at large parameter. Combining the two estimates gives the
required strict bound on $L_y$.

\begin{lemma}[Tangent density of the chord transform]\label{lem:chord-density}
For $\nu$-almost every center $y$,
\begin{equation}\label{chord:transform-density}
 \liminf_{a\to\infty}a^{n/2}J_y(a)
                         \ge\frac{s_n}{2^nW}.
\end{equation}
Consequently,
\begin{equation}\label{eq:chord-upper}
 L_y\le2\left(\frac{W}{s_n}\right)^{2/n}<2
                    \qquad\text{for $\nu$-almost every }y.
\end{equation}
\end{lemma}

\begin{proof}
Use the disjoint chart representation of
Proposition~\ref{prop:chart-mass}. For one chart
$\phi_i:E_i\to Z_i\subset Y$, write
\[
 u_i=v\circ\phi_i,
 \qquad \rho_i=\abs{\theta_i}J_i,
 \qquad J_i=\sqrt{\det G_i},
\]
and extend $u_i$ and $\rho_i$ by zero to $\R^n$. The
bi-Lipschitz bounds give a positive lower bound for $J_i$ on
this chart almost everywhere. Since $\abs{\theta_i}\ge1$,
$v\in L^p(\mu)$ and $\mu(Y)<\infty$, $u_i$ is Lebesgue
integrable. The function $\rho_i$ is integrable as well, by
the mass representation. Thus ordinary Lebesgue differentiation
applies to both functions and to $1_{E_i}$.

For $\nu$-almost every $y=\phi_i(z)$, we may therefore require
all of the following: $z$ is a density point of $E_i$; the centered
metric differential is the Euclidean norm $\abs{\cdot}_{G_0}$,
with $G_0=G_i(z)>0$; and $u_i,\rho_i$ have Lebesgue values
\begin{equation}\label{chord:lebesgue-values}
 v_0=v(y)>0,
 \qquad \rho_0=\abs{\theta_i(z)}\sqrt{\det G_0}>0.
\end{equation}
Indeed the exceptional coordinate sets are Lebesgue-null and
therefore have zero chart mass, while $\{v=0\}$ has zero
$\nu$-mass. There are only countably many charts. The representative
of $v$ agrees with its Lebesgue value at almost every such point.
Fix a center with these properties.

Let $\ell=\liminf_{a\to\infty}a^{n/2}J_y(a)$. If
$\ell=\infty$, the desired lower bound is automatic. Otherwise,
choose $a_j\to\infty$ realizing this lower limit and put
$\varepsilon_j=a_j^{-1/2}$. On every bounded set of
$h\in\R^n$, the functions
\[
 u_i(z+\varepsilon_jh),\qquad
 \rho_i(z+\varepsilon_jh),\qquad
 1_{E_i}(z+\varepsilon_jh)
\]
converge in measure to $v_0,\rho_0,1$, respectively. For example,
this follows by changing variables in the defining mean convergence
at the Lebesgue point $z$. A diagonal subsequence over bounded
balls makes all three convergences hold almost everywhere in
$\R^n$. The subsequence still realizes $\ell$.

At almost every $h$, the points $z+\varepsilon_jh$ consequently
belong to $E_i$ for all sufficiently large $j$. The centered
metric differential then gives
\begin{equation}\label{chord:rescaled-distance}
 R_j(h):=\varepsilon_j^{-1}
          d\bigl(y,\phi_i(z+\varepsilon_jh)\bigr)
                    \longrightarrow\abs{h}_{G_0}.
\end{equation}
When the chart point is absent, define the integrand below to
be zero. Restricting the nonnegative integral defining $J_y$
to this chart and changing variables gives, for every fixed
bounded ball $D_R\subset\R^n$ centered at zero,
\[
 W a_j^{n/2}J_y(a_j)
 \ge\int_{D_R}
 \frac{1_{E_i}(z+\varepsilon_jh)
        u_i(z+\varepsilon_jh)^p\rho_i(z+\varepsilon_jh)}
      {\bigl(1+\kappa^2R_j(h)^2
                   u_i(z+\varepsilon_jh)^{2\beta}\bigr)^n}\dd h.
\]
The factor $a_j^{n/2}$ has canceled the Jacobian
$\varepsilon_j^n$. Fatou's lemma,
\eqref{chord:lebesgue-values}, and~\eqref{chord:rescaled-distance}
now imply
\[
 W\ell\ge v_0^p\rho_0
        \int_{D_R}\bigl(1+v_0^{4\beta}\abs{h}_{G_0}^2\bigr)^{-n}
                 \dd h.
\]
This step uses no uniform integrability of the rescaled densities.
Letting $R\to\infty$ and applying monotone convergence gives the
same bound with $D_R$ replaced by $\R^n$.

The linear substitution $H=v_0^{2\beta}G_0^{1/2}h$ has Jacobian
$v_0^{2n\beta}\sqrt{\det G_0}$. Since $2n\beta=p$, we obtain
\begin{align*}
 W\ell
 &\ge \abs{\theta_i(z)}
                \int_{\R^n}(1+\abs{H}^2)^{-n}\dd H\\
 &=\abs{\theta_i(z)}\,2^{-n}s_n
 \ge2^{-n}s_n.
\end{align*}
The middle equality is the stereographic integral
\eqref{eq:stereographic-integral}. The final inequality uses the
nonzero integer multiplicity $\abs{\theta_i(z)}\ge1$. This proves
\eqref{chord:transform-density} without a bound on the multiplicity.

If $L_y>0$, Proposition~\ref{prop:chord-transform} and
\eqref{chord:transform-density} give
\[
 \frac{s_n}{2^nW}\le(2L_y)^{-n/2}.
\]
Rearranging yields~\eqref{eq:chord-upper}. If $L_y=0$, that
upper bound is immediate. Its strict inequality follows from
$W<s_n$.
\end{proof}

\begin{remark}[The round sphere]\label{rem:chord-sphere}
The constants can be checked simultaneously in every dimension
on the unit round $S^n\subset\R^{n+1}$ with $v=1$. Here
$W=s_n$, $\nu$ is normalized area, and $L_y=2$ by
rotational symmetry. In stereographic coordinates chosen so that
$\abs{x-y}^2=4\abs{z}^2/(1+\abs{z}^2)$,
\[
 J_y(a)=\frac{2^n}{s_n}
     \int_{\R^n}\bigl(1+(1+4a)\abs{z}^2\bigr)^{-n}\dd z
          =(1+4a)^{-n/2}.
\]
Thus the transform estimate is an equality, and its limiting
coefficient is $2^{-n}$, exactly the value in
\eqref{chord:transform-density} when $W=s_n$.
\end{remark}

Only the scalar values $v(y),v(x)$ and the distance $d(y,x)$
occur in~\eqref{eq:chord-upper}. The kernel
\[
 (y,x)\longmapsto
          v(y)^{2\beta}d(y,x)^2v(x)^{2\beta}
\]
is Borel and is integrable against $\nu\otimes\nu$, since $Y$
has finite diameter and $v^{p+2\beta}\in L^1(\mu)$. Accordingly,
the almost-everywhere upper bound can be averaged without any
jointly measurable choice of the covectors $D_xd(y,x)$.
